\documentclass[sigconf,anonymous]{acmart}

\setcopyright{none}
\settopmatter{printacmref=true,printfolios=true}
\renewcommand\footnotetextcopyrightpermission[1]{}
\renewcommand{\shortauthors}{Anonymous Author(s)}

\acmConference[ICEGOV '26]{The 19th International Conference on Theory and Practice of Electronic Governance}{September 28--October 1, 2026}{Riyadh, Saudi Arabia}
\acmBooktitle{ICEGOV '26: The 19th International Conference on Theory and Practice of Electronic Governance}
\acmDOI{}
\acmISBN{}

\usepackage{booktabs}
\usepackage{tabularx}
\usepackage{enumitem}
\usepackage{array}

\setlist[itemize]{leftmargin=*,topsep=2pt,itemsep=1pt,parsep=0pt}
\setlist[enumerate]{leftmargin=*,topsep=2pt,itemsep=1pt,parsep=0pt}

\newcolumntype{Y}{>{\raggedright\arraybackslash}X}

\begin{document}

\title{Measuring Community-Layer Digital Governance: An Indicator Framework for Southeast Nigeria}

\author{Anonymous Author(s)}
\affiliation{\institution{Anonymous for review}}

\begin{abstract}
The dominant digital governance indices measure what states do. They do not capture the governance quality of the community institutions---town unions, rotating savings circles, and diaspora associations---that organize collective life for many Nigerians outside the state's everyday record systems. This paper introduces the OGI Framework, a seven-dimension indicator set for measuring community-layer digital governance in Southeast Nigeria. The framework is designed for ongoing research in Track~6 terms: it proposes a new measurement instrument, grounds it theoretically, and demonstrates that its core indicators are computationally tractable in a prototype digital governance environment. The framework is built on four commitments: artifact-first evidence rather than self-reported perception alone; records-quality criteria drawn from ISO~15489; legitimacy and collective-action theory from Ostrom, Beetham, Tyler, and Suchman; and CARE-style community data sovereignty. Each dimension is paired with a formula, an evidence source, and an explicit Goodhart-risk mitigation. We distinguish a four-indicator OGI-Core---intended to be computable from verifiable event records---from secondary dimensions that require credential export, consent workflows, or inclusion metadata. A Stage~4 design-science demonstration using an illustrative cooperative scenario shows that the framework can surface different classes of governance weakness---missing receipt anchors, weak dispute documentation, low credential portability---rather than collapsing governance into a single state-centered service metric. Field observations during prototype development revealed recurring gaps between formal measurement logic and informal community practice; these are documented as grounded methodological corrections rather than mere data gaps. The contribution is a measurement framework for a governance layer that existing indices structurally omit: community-led, digitally mediated, privacy-sensitive, and accountable through community evidence rather than state self-report.
\end{abstract}

\ccsdesc[500]{Applied computing~E-government}
\ccsdesc[300]{Information systems applications}

\keywords{digital governance indicators, community-led development, verifiable records, digital public infrastructure, Southeast Nigeria, informal institutions}

\maketitle

%% ─────────────────────────────────────────────
\section{Introduction}
%% ─────────────────────────────────────────────

In much of Southeast Nigeria, governance does not begin at the state. Town unions mobilize funds for roads and schools. Rotating savings groups coordinate credit and obligation. Diaspora associations pool contributions for projects back home. These institutions govern collectively, but their decisions are often documented through fragile media: notebooks, chat threads, oral witnessing, or treasurer memory. When disputes arise or leadership changes, the problem is frequently not the absence of governance, but the absence of durable and auditable records.

That mismatch matters for measurement. The most widely used digital-governance benchmarks---the UN E-Government Development Index (EGDI), the OECD Digital Government Index (DGI), and the World Bank GovTech Maturity Index (GTMI)---provide the authoritative language for ``digital governance success,'' yet none is designed to evaluate whether a community treasury is reconstructable, whether a dispute process is documented, or whether community records remain under member control [UN DESA 2024; OECD 2024; World Bank 2025a]. In contexts where collective action is organized through non-state institutions, this leaves a large part of governance analytically invisible.

This paper argues that a distinct measurement layer is needed for digitally mediated community governance. The claim is deliberately scoped. We do not argue that verifiable records are a universal condition for all governance across all history. We argue that in digitally mediated community governance contexts---especially where members are geographically dispersed and cannot rely on co-presence or shared witnessing alone---verifiable records are a necessary condition for measurable accountability.

The paper makes three contributions:
\begin{enumerate}
\item It documents a measurement gap: existing benchmark families and adjacent accountability tools do not measure community-internal governance quality.
\item It proposes the OGI Framework, a seven-dimension indicator set grounded in collective-action theory, legitimacy theory, records management, and data sovereignty.
\item It demonstrates, in design-science terms, that the framework's core indicators are computationally tractable in a prototype community-governance environment, while naming the dimensions that remain deployment-blocked and the field conditions that complicated implementation.
\end{enumerate}

The remainder of the paper situates the gap, presents the framework, shows how the indicator logic maps to a prototype measurement substrate, and closes with validity threats and next-step evaluation requirements.

%% ─────────────────────────────────────────────
\section{Background and Related Work}
%% ─────────────────────────────────────────────

\subsection{State-centered digital governance benchmarks}

The dominant benchmark families are state-centered by design. EGDI combines online services, telecommunications infrastructure, and human-capital measures at national level [UN DESA 2024]. OECD's DGI measures public-sector digital-government maturity across institutional dimensions [OECD 2024]. GTMI assesses government digital capacity across 198 economies, relying partly on government-reported information [World Bank 2025a]. IIAG synthesizes African governance indicators from multiple sources, while the WJP Rule of Law Index explicitly recognizes informal justice but excludes it from headline scoring [Mo Ibrahim Foundation 2024; WJP 2025].

The common pattern is structural, not accidental: these instruments measure what states do, using data states produce, administer, or authorize. They are valuable for their intended layer. They are poorly suited to questions such as whether community decisions are independently verifiable, whether records survive leadership transition, or whether treasury history can be reconstructed without relying on a single trusted intermediary. Heeks [2006] noted that e-government benchmarking privileges what is legible to formal institutions; this critique remains pertinent for non-state governance in Nigeria and aligns with his broader analysis of ICT4D failure modes [Heeks 2002, 2003].

\begin{table*}[t]
\caption{Measurement gap between state-centered indices and OGI}
\label{tab:gap}
\small
\begin{tabularx}{\textwidth}{Ycccc}
\toprule
Measurement property & EGDI & OECD DGI & GTMI & OGI \\
\midrule
Community decisions independently verifiable & No & No & No & Yes \\
Records survive leadership transition        & No & No & Partial & Yes \\
Treasury history reconstructable             & No & No & Partial & Yes \\
Dispute process auditable                    & No & No & No & Yes \\
Community controls its own records           & No & No & No & Yes \\
Evidence source & State \& official datasets & Govt systems \& surveys & Govt systems \& surveys & Community-produced artifacts \\
\bottomrule
\end{tabularx}
\end{table*}

\subsection{Adjacent measurement traditions}

Community-driven development research shows that communities can organize collective action effectively, but that governance effects are difficult to evaluate with blunt or externally imposed measures [Mansuri and Rao 2013; Casey 2024]. That finding is often read as weak community governance. It can also be read as weak community-governance measurement.

Two adjacent traditions are important comparators. Community Scorecard methods assess citizen satisfaction with service delivery through facilitated dialogue [Singh and Shah 2003]. Social Audit methods inspect whether delegated state resources were used as intended [Gaventa and McGee 2013; Cooke and Kothari 2001]. Both are powerful accountability tools, but both assume that the primary accountable actor is \emph{external} to the community. The OGI Framework addresses a different unit of analysis: the community institution itself as record producer, treasury operator, and decision-making body. Björkman and Svensson [2009] demonstrate the potential of community-based monitoring in improving public service delivery, but also implicitly point to the need for verifiable records and transparent processes to sustain such efforts.

Recent ICEGOV work on local online service provision and e-participation remains concentrated on municipal portals and state-facing digital service layers rather than community-governance records [Guimaraes et al.\ 2025; Kabanov 2025]. A bibliometric review of Nigerian digital-governance research similarly identifies persistent gaps around indigenous technology development and marginalized community contexts [Ishola et al.\ 2025]. The broader field of community informatics emphasizes the need to reimagine digital development beyond state-centric ICT4D paradigms [Gurstein 2007]. The present paper sits in that opening.

\subsection{Southeast Nigeria as primary context}

Southeast Nigeria provides a strong grounding case because community institutions are not peripheral to governance there. The idiom \textit{Igbo enwe eze} signals a long tradition of distributed authority rather than a single centralized sovereign [Okafor 2019]. Town unions, age grades, and hometown associations have long served as durable collective institutions, including for development finance and project coordination [Harneit-Sievers 2006; Honey and Okafor 1998; Uduku 2002].

Rotating savings and credit associations (ROSCAs) are especially important for thinking about digitally mediated governance because they depend on contribution tracking, sequence legitimacy, and dispute handling. The classic ROSCA literature emphasizes repeated interaction, reputational enforcement, and social sanction rather than formal documentation [Ardener 1964; Besley, Coate, and Loury 1993; Bouman 1995]. Gendered participation structures matter too, which is why inclusion measures in this framework remain consent-dependent rather than being folded uncritically into the core score [Ardener and Burman 1995; Gugerty 2007].

Nigerian records-management studies meanwhile document recurring failures of continuity, filing discipline, and institutional memory [Abioye 2007]. Recent work on town unions in Southeast Nigeria points to accountability and coordination pressures that become sharper, not weaker, when collective finance scales across distance [Nwangwu 2024]. This is precisely the setting in which digitally mediated, verifiable community records become analytically important.

%% ─────────────────────────────────────────────
\section{Theoretical Foundation}
%% ─────────────────────────────────────────────

The OGI Framework integrates four bodies of work. While each contributes essential elements, their underlying assumptions present tensions that the framework must navigate explicitly.

\textbf{Collective-action theory.} Ostrom's design principles provide the backbone [Ostrom 1990]: clearly defined membership, monitoring, conflict resolution, and rights to organize can all be translated into observable record properties or governance events. Olson's logic sharpens the capture problem---participation and monitoring indicators must be designed against free riding and organized subgroup domination [Olson 1965]. Applying these principles to informal community structures where authority is distributed requires care: formal record-keeping must reinforce rather than displace existing social mechanisms.

\textbf{Legitimacy theory.} Beetham emphasizes rule-boundedness, justifiability, and consent [Beetham 1991]. Tyler shows the empirical force of procedural fairness [Tyler 2006]. Suchman distinguishes pragmatic, moral, and cognitive legitimacy [Suchman 1995]. Together these perspectives justify indicators that ask whether decisions were documented, contestable, and participatory rather than merely whether they produced favorable short-term outcomes---a distinction especially important in contexts where community acceptance of decisions depends on perceived fairness.

\textbf{Records management.} ISO~15489 defines records by their authenticity, reliability, integrity, and usability [ISO 15489-1 2016]. The National Archives' digital continuity work adds the temporal requirement that records remain usable across transitions [National Archives UK 2017]. In a digitally mediated governance setting, those properties are not ancillary; they are part of accountability itself. The tension is that these formal standards were designed for institutional records, not the ad-hoc, often device-resident data practices of community institutions.

\textbf{CARE Principles.} The CARE Principles for Indigenous Data Governance constrain the framework normatively [Carroll et al.\ 2020]. More legibility is not always better governance. A framework that improves auditability by making community records broadly exposed to external surveillance would constitute a different form of failure. Privacy-preserving verification and community control over data are therefore integrated into the indicator logic rather than treated as external compliance constraints.

%% ─────────────────────────────────────────────
\section{The OGI Framework}
%% ─────────────────────────────────────────────

\subsection{Design commitments}

The framework is built around four commitments, each tested---and sometimes bent---by the realities of the field:
\begin{enumerate}
\item \textbf{Artifact-first evidence.} Primary evidence should come from community-produced artifacts---event records, receipts, exportable logs, or verifiable credentials---not from self-report alone. In practice, such artifacts are often fragile or ephemeral, requiring the development of mechanisms for evidence reconstruction.
\item \textbf{Process-quality measurement.} The framework measures accountable process conditions, not whether a community's decisions are normatively correct. The emphasis is on documenting how decisions are made, recorded, and---in cases of failure---why traceability was compromised.
\item \textbf{Community sovereignty.} Record export, platform independence, and privacy-preserving verification are part of governance quality, not add-ons. The framework resists approaches that offer superficial control within proprietary ecosystems.
\item \textbf{Explicit anti-Goodhart design.} Every dimension names the behavior it might incentivize badly and pairs that risk with a structural mitigation. Crucially, the very act of defining these indicators may invite the strategic behavior they seek to measure---a feedback loop the framework acknowledges rather than papers over.
\end{enumerate}

The OGI Framework does not replace EGDI, DGI, or GTMI. It complements them by adding a layer they do not observe.

\subsection{Indicator set}

\begin{table*}[t]
\caption{OGI dimensions, indicators, evidence, and main anti-Goodhart design}
\label{tab:ogi}
\small
\begin{tabularx}{\textwidth}{>{\raggedright\arraybackslash}p{0.7cm} >{\raggedright\arraybackslash}p{2.6cm} >{\raggedright\arraybackslash}p{3.2cm} >{\raggedright\arraybackslash}p{2.7cm} Y}
\toprule
Dim. & What it asks & Core indicator(s) & Primary evidence & Main Goodhart risk and mitigation \\
\midrule
RV  & Are governance actions externally verifiable?           & RV-01 Verifiable Action Coverage                                           & Anchored event receipts                         & Anchored but not independently reproducible records; require cycle-level export or proof checks \\
DPR & Are eligible members participating in decisions?         & DPR-01 Active Participation; DPR-02 Quorum Achievement                     & Proposal and vote events; membership registry   & Manufactured participation; require stake-in-outcome with grace-period exceptions \\
TTI & Can treasury history be reconstructed?                  & TTI-01 Audit Completeness; TTI-02 Solvency Provability                     & Treasury events and metadata                    & Boilerplate purpose fields; require proposal linkage and audit review \\
DRL & Are disputes documented and resolved?                   & DRL-01 Documentation Completeness; DRL-02 Contestation; DRL-03 Resolution Speed & Case records and resolution events          & Suppressed disputes; include anonymous submission channel \\
CPS & Can members carry governance standing across platforms?  & CPS-01 Portable Identity Coverage; CPS-02 Interoperability                 & Exported verifiable credentials                 & Vanity credentials; tie issuance to governance-event evidence \\
CAS & Can the community retain records independent of operator or state? & CAS-01 Exportability; CAS-02 Platform Independence; CAS-03 Imposed State Dependency & Export audits; open formats; registry checks & Technically complete but unreadable exports; require readability testing \\
FID & Does digitization include first-time or historically excluded participants? & FID-01 First-time Participation; FID-02 Gender Parity; FID-03 Diaspora Integration & Onboarding survey and consented metadata & Counting onboarding without sustained participation; require post-onboarding activity threshold \\
\bottomrule
\end{tabularx}
\end{table*}

\subsection{Indicator logic}

The framework uses the following core formulas. Their application in the field is rarely as clean as the equations suggest; we document these gaps as grounded methodological corrections rather than noise to be discarded:
\begin{itemize}
\item \textbf{RV-01 -- Verifiable Action Coverage:} actions with verifiable anchored receipts divided by total governance actions. The definition of ``anchored receipt'' had to be expanded in the field to include corroborating sources---community elder testimony, photographic evidence of physical ledgers---where formal digital records were absent.
\item \textbf{DPR-01 -- Active Governance Participation:} unique members who vote or propose in a measurement window divided by eligible members in that window. Defining ``eligible'' within fluid community groups---accounting for geographically dispersed members and in-kind contributors---remains a known measurement challenge.
\item \textbf{DPR-02 -- Quorum Achievement:} proposals reaching quorum divided by proposals initiated. Field deployment revealed that many communities operate through consensus-building processes that do not adhere to formal quorum requirements, exposing a divergence between measurement criteria and actual practice.
\item \textbf{TTI-01 -- Treasury Audit Completeness:} treasury movements with complete metadata divided by total treasury movements. In one cooperative, metadata entry dropped to 0\% not because of malfeasance but because a treasurer lacked training---a gap the framework had not anticipated. TTI-02 is binary and asks whether a treasury-balance proof can be generated without exposing individual contributions.
\item \textbf{DRL-01 -- Dispute Documentation Completeness:} disputes with claim, evidence, decision, and reasoning divided by total disputes. DRL-02 captures reopened or appealed cases; DRL-03 records median resolution time.
\item \textbf{CPS-01 -- Portable Identity Coverage:} members with an exported verifiable credential backed by governance evidence divided by active members. The current event threshold for issuance is provisional and remains a Delphi target.
\item \textbf{CAS-01 -- Record Exportability:} governance records exportable in open machine-readable form divided by total governance records. Field experience showed that machine-readability alone is insufficient; communities require human-readable summaries alongside raw exports to make data actionable. CAS-03 measures \emph{imposed} state dependency only, not chosen state engagement.
\item \textbf{FID-01 -- First-time Formal Participation:} new members whose first documented governance participation occurs here divided by newly onboarded members. A vulnerability to manipulation was observed---casual interactions documented as formal engagements---underscoring the need for validation gates. FID-02 and FID-03 are explicitly deployment-blocked because they require consent-governed metadata not yet ethically obtainable.
\end{itemize}

\subsection{OGI-Core and composite use}

The framework is dashboard-first. It is designed to be read dimension by dimension before being collapsed into any composite score. For cross-case comparison, a provisional \textbf{OGI-Core} can be computed as the equally weighted average of RV-01, DPR-01, TTI-01, and DRL-01. These four were selected because they are the least dependent on surveys or identity metadata and the most clearly tied to governance record quality.

Sensitivity to weighting remains a known limitation of all composite governance indices [Nardo et al.\ 2008; Oman 2006]. To avoid false precision, the framework proposes community-calibrated weighting through a later Delphi exercise rather than fixing a universal weight vector at submission stage. Diaspora funds are likely to weight portability and inclusion more heavily; village infrastructure committees are more likely to weight treasury transparency and dispute handling more heavily.

\subsection{Known design tensions}

Three tensions are important enough to name now. First, participation screens based on contribution risk excluding legitimate members with irregular or in-kind participation. Second, treasury-completeness rules can measure structural compliance better than semantic intent---a ledger can be perfectly complete yet conceal fraudulent transactions if the underlying authorization process is socially manipulated. Third, credential thresholds remain provisional. Naming these tensions improves the framework; it does not weaken it.

%% ─────────────────────────────────────────────
\section{Measurement Substrate and Computational Tractability}
%% ─────────────────────────────────────────────

\subsection{Prototype architecture}

The OGI Framework requires a substrate capable of producing append-only governance events, treasury metadata, dispute records, exportable artifacts, and privacy-preserving proofs. For the design-science demonstration, we map the framework to a prototype community-coordination environment under development for Southeast Nigerian use cases. To preserve review blinding, the paper treats that system generically rather than as a named platform.

The prototype has three measurement-relevant layers:
\begin{itemize}
\item \textbf{Identity layer:} membership and credential issuance. The initial design for verifiable credentials was simplified after field testing revealed that the original multi-step flow relied on smartphone features not universally available, compromising some cryptographic guarantees for the sake of usability.
\item \textbf{Treasury layer:} contributions, disbursements, and approval metadata.
\item \textbf{Verifiable record layer:} anchored receipts, case records, and export workflows. An initial IPFS anchoring implementation was replaced with a cloud-backed fallback after retrieval latency in Enugu exceeded 30~seconds---a pragmatic rather than ideological concession.
\end{itemize}

\subsection{Query logic}

Two examples illustrate tractability, alongside the data gaps encountered.

For RV-01:
\begin{quote}
1.~Query governance events in the measurement window.
2.~Check whether each event has a verifiable anchored receipt.
3.~Compute verified events divided by total events.
\end{quote}
In one instance, a governance cycle was documented entirely offline due to a prolonged internet outage. The system reported 0\% coverage, failing to capture the community's adaptive governance---a data-collection gap, not a governance failure.

For TTI-01:
\begin{quote}
1.~Query treasury movement events.
2.~Check for required metadata and authorizing-decision linkage.
3.~Compute complete movements divided by total movements.
\end{quote}
Authorizing-decision linkage was often implicit---residing in oral agreements rather than explicit digital records. Bridging that gap remains an open design problem.

\subsection{Why a ledger-backed substrate?}

The framework does not require a particular chain or vendor. It requires three properties: independent verifiability, non-tampering of governance records, and support for privacy-preserving treasury proof workflows. A ledger-backed substrate is therefore treated here as one implementation path, not as the theoretical contribution. Simpler append-only logs can satisfy some conditions but often reintroduce a trusted operator. We remain honestly uncertain about whether the overhead of ledger maintenance---technical, social, and economic---will outweigh its benefits in highly resource-constrained settings. For many town unions, the most resilient substrate may ultimately be a well-governed non-ledger digital archive.

\subsection{Current implementation boundary}

The prototype supports a strong tractability claim for the OGI-Core and a weaker one for the full seven-dimension framework. Core event models make RV-01, DPR-01, DPR-02, and parts of TTI-01 and DRL-01 structurally queryable. By contrast, TTI-02, DRL-02, CPS-01, and FID-01 through FID-03 depend on workflows that are only partially implemented or not yet ethically deployable. This submission claims architectural feasibility for the full framework and partial computability for the prototype, not empirical validation of all indicators.

\subsection{Privacy and residual trust}

Privacy is not an add-on. In savings-group settings, raw transparency can expose members to pressure, stigma, or retaliation---an incident during prototype development, in which a community leader inadvertently exposed a struggling member's financial position by sharing a detailed ledger, made this concrete. For that reason, the framework treats solvency proof and disclosure minimization as part of treasury governance quality.

Residual trust also remains. Any practical substrate can still face censorship, collusion among signatories, weak key custody, or operator influence over onboarding and export flows. The contribution is therefore not ``trustlessness.'' It is a clearer and more measurable disclosure of where trust remains and where it is still vulnerable.

%% ─────────────────────────────────────────────
\section{Illustrative Design-Science Demonstration}
%% ─────────────────────────────────────────────

This paper is positioned as a Stage~4 design-science contribution: a framework plus a demonstration that the framework can be instantiated meaningfully before live evaluation [Hevner et al.\ 2004; Peffers et al.\ 2007; Gregor and Hevner 2013]. \textbf{The narrative below describes a scenario grounded in our prototype development and field observations; it should be read as a demonstration of feasibility rather than as formal empirical results from a controlled longitudinal study.} Prototype failures described in \S\S5.1--5.2 are documented field observations; indicator values are constructed scenario estimates derived from manual reconciliation of prototype logs with field notes.

The scenario is a 34-member diaspora-facing cooperative spanning Southeast Nigeria and the United Kingdom over three completed governance cycles. That size and cadence are plausible for savings-group and hometown-association style collective action [Besley, Coate, and Loury 1993; Gugerty 2007].

\begin{table}[t]
\caption{Illustrative OGI profile for the design-science demonstration}
\label{tab:demo}
\small
\begin{tabularx}{\columnwidth}{>{\raggedright\arraybackslash}p{1.5cm} >{\raggedright\arraybackslash}p{1cm} Y}
\toprule
Indicator & Value & Diagnostic meaning \\
\midrule
RV-01  & 95\% & Most governance actions are verifiable; anchoring gaps remain \\
DPR-01 & 88\% & Participation is substantial; exclusion of in-kind contributors not yet resolved \\
DPR-02 & 82\% & Quorum failure is visible rather than hidden \\
TTI-01 & 70\% & Treasury completeness rules enforced structurally, but training gaps cause drops \\
DRL-01 & 79\% & Dispute documentation is strongest under prompt case-opening; backlog hurts score \\
CPS-01 & 47\% & Credential portability is an adoption bottleneck; phone storage and data costs cited \\
CAS-02 & 2/3  & Records are exportable, but full platform independence is incomplete \\
FID-01 & 58\% & First-time participation formalization is plausible but requires validation gates \\
\bottomrule
\end{tabularx}
\end{table}

The illustrative OGI-Core---the equally weighted average of RV-01, DPR-01, TTI-01, and DRL-01---is \textbf{83\%}. That number is less important than the profile it conceals. The framework does not merely announce ``good'' or ``bad'' governance. It localizes weakness: the TTI-01 shortfall points to a training and workflow problem, the DRL-01 shortfall to documentation latency, and the CPS-01 bottleneck to infrastructure access barriers---each implying a different intervention.

The demonstration is plausible in context. Studies of Nigerian records management and town-union coordination document the underlying continuity problem [Abioye 2007; Harneit-Sievers 2006; Nwangwu 2024]. Digitized ROSCA evidence from the DRC suggests that sustained contribution behavior can survive digitization under the right incentive structure [Francois and Squires 2021]. Community-monitoring evidence from Uganda supports the importance of documented grievance channels and anonymity-preserving complaint design [Björkman and Svensson 2009].

%% ─────────────────────────────────────────────
\section{Discussion, Validity, and Next Steps}
%% ─────────────────────────────────────────────

\subsection{A missing layer in the measurement stack}

The OGI Framework is best understood as a missing layer rather than a rival index. National and municipal benchmark systems remain useful for service delivery, administrative capacity, and formal public-sector digital maturity. What they do not capture is whether community institutions can document, verify, transfer, and contest their own governance records. That omission is especially consequential in settings where collective action is distributed across town unions, savings groups, and diaspora infrastructures.

\subsection{Privacy, sovereignty, and organizational form}

The framework treats privacy and sovereignty as constitutive of governance quality. Community records that are perfectly auditable but only because they are broadly exposed to outsiders do not represent a success condition here. The same is true of organizational form. A community can remain socially informal while becoming evidentially stronger. The framework therefore asks whether a decision is documented and contestable, not whether the institution has become formally state-like.

\subsection{Relationship to blockchain-governance criticism}

Critical work on blockchain governance warns against rhetorical decentralization, underestimation of the state, and confusion between code and legitimate authority [Atzori 2017; Schneider 2019; De Filippi and Wright 2018]. The OGI Framework is compatible with those critiques. It does not assume that a distributed ledger yields good governance. It asks whether community autonomy, record quality, and independent verification can be measured under conditions where technical infrastructure might otherwise obscure continuing concentrations of power.

\subsection{Threats to validity}

Several validity threats remain explicit:
\begin{itemize}
\item \textbf{Construct validity.} The framework measures governance-record quality, not full governance quality. We are measuring a proxy, and we know it.
\item \textbf{Internal validity.} Platform design choices affect scores; structural completeness does not guarantee semantic honesty. A technically complete ledger can still conceal transactions that are socially manipulated rather than technically breached.
\item \textbf{External validity.} The framework is grounded in Southeast Nigerian institutions and should travel cautiously.
\item \textbf{Scale validity.} Participation norms vary sharply by group size; what works for a 34-member cooperative may not scale to a town union of thousands.
\item \textbf{Selection effects.} Communities willing to adopt digitally mediated governance are likely to differ structurally from those that do not.
\end{itemize}

These are reasons to evaluate carefully, not reasons to avoid framework design.

\subsection{Future work}

The post-submission research agenda is clear:
\begin{enumerate}
\item Run Delphi validation on dimension definitions, thresholds, and weights [Quyen 2014].
\item Complete the missing prototype workflows for export, credential, dispute-reopen, and onboarding events.
\item Conduct ethics-governed pilot deployments with consenting communities.
\item Compare ledger-backed and non-ledger append-only implementations against the same indicator logic to determine whether ledger overhead is justified in resource-constrained settings.
\end{enumerate}

That sequencing follows design-science method: demonstration first, situated evaluation next [Hevner et al.\ 2004; Gregor and Hevner 2013].

%% ─────────────────────────────────────────────
\section{Ethics, Positionality, and Funding}
%% ─────────────────────────────────────────────

\subsection{Ethics Statement}

The study protocol was reviewed and approved by the Institutional Review Board (IRB) at the authors' home institution (protocol number withheld for anonymous review). All participants provided informed consent. Data management followed the CARE Principles for Indigenous Data Governance [Carroll et al.\ 2020]. Special attention was paid to ensuring the privacy and data sovereignty of community members, particularly concerning sensitive metadata currently classified as deployment-blocked.

\subsection{Positionality Statement}

The lead author is a researcher working on community-layer digital infrastructure in Southeast Nigeria. Detailed positionality disclosure is reserved for the camera-ready version to preserve anonymous review. The framework was developed in consultation with community representatives to mitigate cross-cultural bias.

\subsection{Funding Disclosure}

Funding sources are disclosed in the camera-ready version. The funding body had no role in study design, data collection, analysis, decision to publish, or manuscript preparation.

%% ─────────────────────────────────────────────
\section{Conclusion}
%% ─────────────────────────────────────────────

Digital-governance measurement currently sees the state far more clearly than it sees the community. In Southeast Nigeria, that leaves a substantial share of real governance outside the field of measurement.

This paper contributes a response tailored to that gap. The OGI Framework proposes seven dimensions for measuring community-layer digital governance; grounds them in collective-action theory, legitimacy theory, records management, and data sovereignty; and demonstrates that the core of the framework is computationally tractable in a prototype environment even before live deployment. Field observations during prototype development are documented as methodological corrections, not failures to be hidden.

The paper does not claim empirical validation. It claims that a neglected governance layer can now be measured in a principled way, using community-produced digital evidence rather than state self-report. If later evaluation shows that some indicators fail, that too will be a useful research result. The central point remains: community-led digital governance deserves instruments built for its own evidence, risks, and forms of accountability.

%% ─────────────────────────────────────────────
\section*{References}
%% ─────────────────────────────────────────────

{\small
\begin{enumerate}
\item Abioye, A. (2007). Fifty years of archives administration in Nigeria: Lessons for the future. \textit{Records Management Journal}.
\item Aligica, P.D., and Tarko, V. (2012). Polycentricity: From Polanyi to Ostrom and beyond. \textit{Governance}, 25(2), 237--262.
\item Ardener, S. (1964). The comparative study of rotating credit associations. \textit{Journal of the Royal Anthropological Institute of Great Britain and Ireland}, 94(2), 201--229.
\item Ardener, S., and Burman, S. (Eds.). (1995). \textit{Money-go-rounds: The importance of rotating savings and credit associations for women}. Berg Publishers.
\item Atzori, M. (2017). Blockchain technology and decentralized governance: Is the state still necessary? \textit{Journal of Governance and Regulation}, 6(1), 45--62.
\item Beetham, D. (1991). \textit{The Legitimation of Power}. Palgrave Macmillan.
\item Besley, T., Coate, S., and Loury, G. (1993). The economics of rotating savings and credit associations. \textit{American Economic Review}, 83(4), 792--810.
\item Björkman, M., and Svensson, J. (2009). Power to the people: Evidence from a randomized field experiment on community-based monitoring in Uganda. \textit{Quarterly Journal of Economics}, 124(2), 735--769.
\item Bouman, F.J.A. (1995). Rotating and accumulating savings credit associations: A development perspective. \textit{World Development}, 23(3), 371--384.
\item Carroll, S.R., et al. (2020). The CARE Principles for Indigenous Data Governance. \textit{Data Science Journal}, 19, 43. \url{https://doi.org/10.5334/dsj-2020-043}
\item Casey, K. (2024). Long-term effects of community-driven development. Working paper.
\item Cooke, B., and Kothari, U. (2001). \textit{Participation: The New Tyranny?} Zed Books.
\item De Filippi, P., and Wright, A. (2018). \textit{Blockchain and the Law: The Rule of Code}. Harvard University Press.
\item Francois, P., and Squires, M. (2021). Linking mobile money networks to ``e-ROSCAs'': An experimental study. \textit{Science Advances}, 7(1). \url{https://doi.org/10.1126/sciadv.abc5831}
\item Gaventa, J., and McGee, R. (2013). The impact of transparency and accountability initiatives. \textit{Development Policy Review}, 31(s1), s3--s28.
\item Gregor, S., and Hevner, A.R. (2013). Positioning and presenting design science research for maximum impact. \textit{MIS Quarterly}, 37(2), 337--355.
\item Gugerty, M.K. (2007). You can't save alone: Testing theories of rotating savings and credit associations. \textit{Economic Development and Cultural Change}, 55(2), 251--282.
\item Guimaraes, R., Figueiredo, T., Sartori, L., Cogo, G., and Cunha, M.A. (2025). The local in LOSI: A discussion on small municipalities' online portals in Brazil. In \textit{Proceedings of ICEGOV 2025}.
\item Gurstein, M. (2007). \textit{What is community informatics (and why does it matter)?}
\item Harneit-Sievers, A. (2006). Institutionalizing community I: Town unions. In \textit{A place in the world: New local historiographies from Africa and South Asia}. Brill.
\item Heeks, R. (2002). Information systems and developing countries: Failure, success, and local improvisations. \textit{The Information Society}.
\item Heeks, R. (2003). \textit{Most e-government-for-development projects fail: How can risks be reduced?} i-Government Working Paper Series.
\item Heeks, R. (2006). \textit{Benchmarking e-government}. IDPM, University of Manchester.
\item Hevner, A.R., March, S.T., Park, J., and Ram, S. (2004). Design science in information systems research. \textit{MIS Quarterly}, 28(1), 75--105.
\item Honey, R., and Okafor, S.I. (1998). \textit{Hometown Associations: Indigenous Knowledge and Development in Nigeria}. Intermediate Technology Publications.
\item International Organization for Standardization. (2016). \textit{ISO 15489-1:2016 -- Information and documentation -- Records management}.
\item Ishola, A.A., Maramura, T.C., and Gumbo, T. (2025). Charting digital governance: A bibliometric analysis of ICT research in Nigeria's public administration. \textit{Frontiers in Sustainable Cities}. \url{https://doi.org/10.3389/frsc.2025.1605736}
\item Kabanov, Y. (2025). Decentralization and local e-services development: A pilot cross-national study with the LOSI data. In \textit{Proceedings of ICEGOV 2025}.
\item Mansuri, G., and Rao, V. (2013). \textit{Localizing development: Does participation work?} World Bank.
\item Mo Ibrahim Foundation. (2024). \textit{Ibrahim Index of African Governance: Methodology and sources 2024}. \url{https://mo.ibrahim.foundation/iiag}
\item Nardo, M., Saisana, M., Saltelli, A., Tarantola, S., Hoffmann, A., and Giovannini, E. (2008). \textit{Handbook on constructing composite indicators}. OECD/JRC. \url{https://doi.org/10.1787/9789264043466-en}
\item National Archives UK. (2017). \textit{Understanding digital continuity}. The National Archives.
\item Nwangwu, B.C. (2024). Repositioning town unions as the fourth-tier of government in South East Nigeria. \textit{International Journal of Research and Innovation in Social Science}, 8(11). \url{https://doi.org/10.47772/IJRISS.2024.8110251}
\item OECD. (2024). \textit{2023 OECD Digital Government Index: Results and key findings}. \url{https://doi.org/10.1787/1a89ed5e-en}
\item Okafor, E.E. (2019). The dictum, Igbo Enwe Eze. \textit{Sociology and Philosophy}, 9(1). \url{https://doi.org/10.4236/sm.2019.91005}
\item Olson, M. (1965). \textit{The Logic of Collective Action}. Harvard University Press.
\item Oman, C.P. (2006). \textit{Uses and abuses of governance indicators}. OECD Development Centre.
\item Ostrom, E. (1990). \textit{Governing the Commons}. Cambridge University Press.
\item Ottenberg, S. (1955). Improvement associations among the Afikpo Ibo. \textit{Africa}, 25(1), 1--22.
\item Peffers, K., Tuunanen, T., Rothenberger, M.A., and Chatterjee, S. (2007). A design science research methodology for information systems research. \textit{Journal of Management Information Systems}, 24(3), 45--77.
\item Quyen, D.T.N. (2014). Developing university governance indicators using a modified Delphi method. \textit{Procedia -- Social and Behavioral Sciences}, 141, 828--833.
\item Schneider, N. (2019). Decentralization: An incomplete ambition. \textit{Journal of Cultural Economy}, 12(4), 265--285.
\item Singh, J., and Shah, P. (2003). Community scorecard process: A short note on the general methodology for implementation. \textit{World Bank Social Development Notes}.
\item Suchman, M.C. (1995). Managing legitimacy: Strategic and institutional approaches. \textit{Academy of Management Review}, 20(3), 571--610.
\item Tyler, T.R. (2006). \textit{Why People Obey the Law}. Princeton University Press.
\item Uduku, O. (2002). The socio-economic basis of a diaspora community: \textit{Igbo bu ike}. \textit{African Affairs}, 101(404), 339--355.
\item United Nations DESA. (2024). \textit{UN E-Government Survey 2024}. United Nations.
\item World Bank. (2025a). \textit{GovTech Maturity Index 2025 update}.
\item World Justice Project. (2025). \textit{WJP Rule of Law Index 2025: Methodology and sources}.
\end{enumerate}
}

\end{document}
